\GalSourcePageBoundary{090}{L0089}{61}{61}
que deux valeurs différentes; donc $m$ ne peut avoir plus de trois
valeurs; donc enfin $p = 3$. Ainsi, c'est seulement dans ce cas
que le groupe réduit pourra contenir des substitutions de l'ordre~$p$.

Et, en effet, la réduite sera alors du quatrième degré, et, par
conséquent, soluble par radicaux.

Nous savons par là qu'en général, parmi les substitutions de notre
groupe réduit, il ne devra pas se trouver de substitutions de
l'ordre~$p$. Peut-il y en avoir de l'ordre~$p - 1$? C'est ce que je vais
rechercher\GalHistoricalFootnoteMark{1}.

\GalHistoricalFootnoteText{1}{J'ai cherché inutilement dans les papiers de Galois la continuation de ce
qu'on vient de lire. \textup{(A.~Ch.)}}

\par\vspace{1.0em}\noindent\dotfill\par
\vspace{0.4em}\hrule\vspace{2.6em}
\begin{center}\bfseries FIN.\end{center}
